\documentclass[10pt,a4paper]{article}

% ----------------- Paquetes -------------------%
\usepackage[T1]{fontenc}
\usepackage[utf8]{inputenc}
\usepackage[a4paper,margin=2.5cm]{geometry}
\usepackage{mathptmx}             
\usepackage{changepage} 
\usepackage{setspace}
\usepackage[spanish]{babel}
\usepackage[labelfont=bf,labelsep=space]{caption}
\usepackage{amssymb}
\usepackage{amsmath}
\usepackage{ebgaramond}
\usepackage{placeins}
\usepackage{etoolbox}
\usepackage{setspace}
\usepackage{parskip} 
\usepackage{tcolorbox}
\usepackage{needspace}
\usepackage{xparse}
\usepackage{ocgx2}
\usepackage{hyperref}
\usepackage{hypcap}


%%%%%%%%%%%%%%%%%%%%%%%%%%%%%%%%%%%%%%%%%%%%%%%%%%%%%%%%%%%%%%%%
% --- Fija primero tus valores globales "buenos" ---
\setstretch{1.2}                 % Interlineado global
\setlength{\parskip}{0.7\baselineskip}
\setlength{\parindent}{0pt}
\newlength{\GlobalParskip}
\setlength{\GlobalParskip}{\parskip}

\newcounter{eqnum}[subsection]
\renewcommand{\theeqnum}{\arabic{eqnum}}
\newcounter{alphaitem}
\newcounter{numitem}

%% ------------------- Recuadros----------------------%%
\newcommand{\puntosclave}[1]{%
  \begin{tcolorbox}[
    colframe=black,
    title={Puntos clave},
    before upper={\setlength{\parindent}{0.5cm}\setlength{\parskip}{0.5\baselineskip}}
  ]
  #1
  \end{tcolorbox}
}

\newcommand{\modificaciones}[1]{
  \begin{tcolorbox}[
    colback=white,
    colframe=red,
    title={Modificaciones a realizar},
    before upper={\setlength{\parindent}{0.5cm}\setlength{\parskip}{0.5\baselineskip}}
  ]
  #1
  \end{tcolorbox}
}

%% ------------------- Preguntas TEST----------------------%%
% Opciones: radio (single-choice)
\NewDocumentCommand{\OptRadio}{m m m}{%
  \noindent\CheckBox[radio,name=#1,value=#2,width=1.2ex,height=1.2ex]{}%
  \;\textbf{#2.}~#3\par
}

% Opciones: checkbox (multi-choice)
\NewDocumentCommand{\OptCheck}{m m}{%
  \noindent\CheckBox[name=#1,width=1.2ex,height=1.2ex,bordercolor={0 0 0}]{}%
  \;#2\par
}

% Pregunta single-choice (radio)
% Args: 1=id, 2=enunciado, 3=opciones, 4=solución+justificación, 5=imagen (o {}), 6=origen
\NewDocumentCommand{\PreguntaTestSingle}{m m m m m m}{%
  \par\medskip
  \noindent\textbf{#1.}~#2\par\smallskip

    % Imagen (si no es {})
    \IfBlankTF{#5}{}{%
      \begin{center}
        \includegraphics[width=0.75\linewidth]{#5}
      \end{center}
    }

    \begin{Form}
      #3
    \end{Form}

    \par\smallskip
    \switchocg{sol-#1}{\fbox{Mostrar / ocultar solución}}%
    \hfill{\footnotesize #6}\par\smallskip

    \begin{ocg}{Solución}{sol-#1}{0}
      \noindent #4
    \end{ocg}
  \par\medskip
}

% Pregunta multi-choice (checkboxes)
\NewDocumentCommand{\PreguntaTestMulti}{m m m m m m}{%
  \par\medskip
  \noindent\textbf{#1.}~#2\par\smallskip

    % Imagen (si no es {})
    \IfBlankTF{#5}{}{%
      \begin{center}
        \includegraphics[width=0.75\linewidth]{#5}
      \end{center}
    }
    \begin{Form}
      #3
    \end{Form}

    \par\smallskip
    \switchocg{sol-#1}{\fbox{Mostrar / ocultar solución}}%
    \hfill{\footnotesize #6}\par\smallskip

    \begin{ocg}{Solución}{sol-#1}{0}
      \noindent #4
    \end{ocg}
  \par\medskip
}



%% ------------------- Listas----------------------%%
    % --- Lista alfabética ---
    \newenvironment{la}
      {\begin{list}{\alph{alphaitem})}{%
          \usecounter{alphaitem}%
          \setlength{\leftmargin}{1cm}%
          \setlength{\parskip}{3pt}% 
          \setlength{\itemsep}{0pt}% ← espacio entre ítems
          \setlength{\parsep}{0pt}%  ← párrafos dentro de un ítem
      }}
      {\end{list}}
      
    % --- Lista numérica ---
    \newenvironment{lnum}
      {\begin{list}{\arabic{numitem}.}{%
          \usecounter{numitem}%
          \setlength{\leftmargin}{1cm}%
          \setlength{\parskip}{3pt}% 
          \setlength{\itemsep}{0pt}% ← espacio entre ítems
          \setlength{\parsep}{0pt}%  ← párrafos dentro de un ítem
      }}
      {\end{list}}

%% ------------------- Preguntas Test ----------------------%%
      \usepackage{xparse} % si ya lo tienes, no lo repitas

    \NewDocumentEnvironment{test}{m m o}{%
      \begin{tcolorbox}[
        colback=white,
        colframe=black!15,
        boxrule=0.6pt,
        arc=2mm,
        left=2mm,right=2mm,top=2mm,bottom=2mm
      ]
      \centering
      \includegraphics[width=\linewidth]{#1}
      \end{tcolorbox}
    
      \vspace{2mm}
    
      % Caja SOLO para texto (SÍ partible y mantiene márgenes al partir)
      \begin{tcolorbox}[
        enhanced,
        breakable,
        colback=black!3,
        colframe=black!25,
        boxrule=0.6pt,
        arc=2mm,
        left=2mm,right=2mm,top=1.5mm,bottom=1.5mm
      ]
      \textbf{Respuesta:} \textbf{#2}%
      \IfValueT{#3}{\par\smallskip \textbf{Justificación:} #3}
      \end{tcolorbox}
    }{}
      
% ------------------- Imágenes----------------------%
    \graphicspath{{Img/}}% Rutas por defecto 
    % --- Imagen adaptativa ---
    % Registros de dimensión definidos una sola vez
    \newdimen\imgcap
    \newdimen\imgmin
    \newdimen\imgmax
    \newdimen\imgremain
    \newdimen\imgh
    \newdimen\imgneed
    
    \newcommand{\imagenfit}[3]{%
      \addvspace{10pt}% separación antes
      \par\begingroup
        % Parámetros de composición (asignaciones, no \newdimen)
        \imgcap=1\baselineskip            % alto estimado de título+fuente
        \imgmin=0.15\textheight
        \imgmax=0.15\textheight
        \imgremain=\pagegoal
        \ifdim\pagegoal=\maxdimen \imgremain=0pt\fi
        \advance\imgremain by -\pagetotal
        \advance\imgremain by -\imgcap
        % Decisión de altura destino
        \ifdim\imgremain<\imgmin
          \newpage
          \imgh=0.20\textheight
        \else
          \imgh=\imgremain
          \ifdim\imgh>\imgmax \imgh=\imgmax \fi
        \fi
        % Reservar espacio para evitar cortes del bloque completo
        \imgneed=\imgh \advance\imgneed by \imgcap
        \Needspace{\imgneed}
        % Cuerpo
        \noindent\begin{minipage}{\textwidth}
          \centering
          \captionof{figure}{\textemdash\ #2}\par\vspace{0.3em}% título arriba
          \includegraphics[width=\linewidth,height=\imgh,keepaspectratio]{\detokenize{#1}}\par
          \vspace{0.3em}\small\textit{Fuente: }#3% fuente abajo
        \end{minipage}%
      \endgroup
      \par\addvspace{10pt}%
    }

%------------ Bloque de RECUADRO ESPECIAL -------------%
    \newenvironment{specialblock}[1]{%
      \begin{quote} % crea sangría a izquierda y derecha
      \small % texto más pequeño
      \noindent\textbf{#1}\\[-0.3em] % título en negrita
      \rule{\linewidth}{0.4pt}\\[0.6em]}{
      \end{quote}}
      
%-------------------------- Portada -----------------------%
\newcommand{\oldtitle}[1]{\def\OldTitle{#1}}
\newcommand{\changerationale}[1]{\def\ChangeWhy{#1}}

\newcommand{\changebox}{%
\begin{tcolorbox}[title={Historial de cambios del tema}]
\textbf{Título anterior:} \OldTitle\par
\textbf{Motivo del cambio:} \ChangeWhy
\end{tcolorbox}
}

%-------------------------- Author Font -----------------------%
    \newcommand{\authorfont}[1]{{\fontfamily{EBGaramond-TLF}\selectfont #1}}

%-------------------------- Anexos -----------------------%
    \makeatletter
    \newcounter{anexo}
    \renewcommand{\theanexo}{\Roman{anexo}}
    \newcommand{\anexo}[1]{%
      \clearpage
      \phantomsection
      \refstepcounter{anexo}%
      \noindent\textbf{Anexo \theanexo.- }#1\par\medskip
      \addcontentsline{stc}{section}{Anexo \theanexo.- #1}%
    }
    \makeatother

    %-------------------------- Lista de figuras (personal) -----------------------%
\makeatletter
\newcommand{\MyListOfFigures}{%
  \clearpage
  \phantomsection
  {\centering\bfseries\Large Índice de figuras\par}%
  \medskip
  % Entrada en nuestro índice de contenido (stc)
  \addcontentsline{stc}{section}{Índice de figuras}%
  % Recuperación del listado (sin cabecera propia)
  \begingroup
    \parindent=0pt\parskip=2pt
    \@starttoc{lof}%
  \endgroup
}
\makeatother

%-------------------------- Bibliografía -----------------------%
\makeatletter
% Contador para las referencias
\newcounter{myref}
% Cabecera de la bibliografía, entra en el índice 'stc'
\newcommand{\MyBibliography}{%
  \clearpage
  \phantomsection
  {\centering\bfseries\Large Bibliografía y referencias\par}%
  \medskip
  \addcontentsline{stc}{section}{Bibliografía y referencias}%
  \setcounter{myref}{0}%
}
\newcommand{\MyBibItem}[4]{%
  \refstepcounter{myref}%
  \noindent[\themyref]\ #1.\ \emph{#2}. #3. #4\par
}
\makeatother

%--------- Establecer cuatro niveles de índice --------%
    \setcounter{tocdepth}{4}   
    \setcounter{secnumdepth}{4}% numera hasta \paragraph

%%%%%%%%%%%%%%%%%%%%%%%%%%%%%%%%

\makeatletter

    %--------- Interlineado Global --------%
    \edef\GlobalBaseLineStretch{\baselinestretch}
    
    %--------- Espaciado de seccinones --------%
    \renewcommand\section{\@startsection{section}
      {1}{\z@}
      {2.5ex \@plus 1ex \@minus .2ex} % espacio antes
      {1.5ex \@plus .2ex}             % espacio después
      {\normalfont\Large\bfseries}    % formato del título
    }
    \renewcommand\subsection{\@startsection{subsection}{2}{\z@}
      {3ex \@plus 0.8ex \@minus .2ex}
      {1.5ex \@plus .2ex}
      {\normalfont\large\bfseries}
    }
    \renewcommand\subsubsection{\@startsection{subsubsection}{3}{\z@}
      {3ex \@plus 0.5ex \@minus .2ex}
      {1ex \@plus .2ex}
      {\normalfont\normalsize\bfseries}
    }
    \renewcommand\paragraph{\@startsection{paragraph}{4}{\z@}%
    {-2ex \@plus -0.5ex \@minus -.2ex}% espacio antes
    {0.8ex \@plus .2ex}% espacio después
    {\normalfont\normalsize\itshape}}% ← cursiva en lugar de negrita

    %--------- Ecuaciones --------%   
    % Variables globales para almacenar títulos y notas
    \gdef\leq@current@section{}%
    \gdef\leq@current@subsection{}%
    \gdef\leq@last@printed@section{}%
    \gdef\leq@last@printed@subsection{}%
    
    % Comando para añadir nota (invisible en el cuerpo)
    \newcommand{\eqnote}[1]{%
      \addcontentsline{leq}{leqnote}{#1}%
    }
    
    % Comando para ecuaciones
    \newcommand{\eqblock}[2]{%
      % Verificar si necesitamos imprimir el título de sección
      \ifx\leq@current@section\leq@last@printed@section\else
        \addcontentsline{leq}{leqsection}{\leq@current@section}%
        \global\let\leq@last@printed@section\leq@current@section
        \gdef\leq@last@printed@subsection{}% Reset subsección al cambiar sección
      \fi
      % Verificar si necesitamos imprimir el título de subsección
      \ifx\leq@current@subsection\leq@last@printed@subsection\else
        \ifx\leq@current@subsection\@empty\else
          \addcontentsline{leq}{leqsubsection}{\leq@current@subsection}%
          \global\let\leq@last@printed@subsection\leq@current@subsection
        \fi
      \fi
      % Ecuación
      \refstepcounter{eqnum}%
      \begin{equation*}
        #1\tag{E.\theeqnum}
      \end{equation*}%
      \addcontentsline{leq}{leqt}{[E.\theeqnum] -- #2}%
      \addcontentsline{leq}{leqm}{#1}%
    }
    
    %%% ----- Índice de ecuaciones ----------%%%
    % Tamaño para []
    \newcommand{\leqIndexFont}{\normalsize}
    \newcommand{\leqTitleFont}{\tiny}
    \newcommand{\leqNoteFont}{\tiny}
    % Cabecera de sección 
    \newcommand*\l@leqsection[2]{%
      \addvspace{25pt}% Mayor separación antes de nueva sección
      \noindent\leqIndexFont
      \makebox[\linewidth]{\rule{.38\linewidth}{0.4pt}\ \ \textbf{  \MakeUppercase{#1}}\ \ \rule{.38\linewidth}{0.4pt}}%
      \par\vspace{-10pt}
    }
    % Cabecera de subsección 
    \newcommand*\l@leqsubsection[2]{%
      \par\addvspace{15pt}%
      \noindent\leqIndexFont #1
      \par\vspace{-7pt}
    }
    % Título de ecuación 
    \newcommand*\l@leqt[2]{%
    \par\addvspace{10pt}%
      \par\noindent\leqTitleFont #1\par
    }
    % Miniatura de ecuación
    \newcommand*\l@leqm[2]{%
      \vspace{0.3\baselineskip}%
      \par\scriptsize\noindent\hspace*{1.5em}\(#1\)\par%
    }
    % Nota de ecuación 
    \newcommand*\l@leqnote[2]{%
      \vspace{0.2\baselineskip}%
      \par\noindent\leqNoteFont\hspace*{1.5em}\textbf{Nota.--} #1\par\vspace{0.5\baselineskip}%
    }
    % Generación del índice
    \newcommand{\mylistofequations}{%
      \clearpage
      \phantomsection
      % Título visual de la lista (ajústalo a tu gusto):
      {\centering\bfseries\Large Índice de ecuaciones\par}
      % Entrada en nuestro índice 'stc':
      \addcontentsline{stc}{section}{Índice de ecuaciones}%
      % Llamada a tu lista real (usa el nombre exacto de tu macro):
      \@starttoc{leq}%
    }
    
    %--------- Captura de títulos de sección --------%
    \let\leq@original@section\section
    \let\leq@original@subsection\subsection
    % Flag para saber si estamos en una sección asterisco
    \newif\ifLeqStarSection
    \LeqStarSectionfalse
    
    % Redefinir \section CON CONTROL DE ASTERISCO
    \renewcommand{\section}{%
      \@ifstar{\leq@section@star}{\@dblarg\leq@section@nostar}%
    }
    \def\leq@section@star#1{%
      \global\LeqStarSectiontrue
      \gdef\leq@current@section{#1}%
      \gdef\leq@current@subsection{}%
      \leq@original@section*{#1}%
      \global\LeqStarSectionfalse
    }
    \def\leq@section@nostar[#1]#2{%
      \global\LeqStarSectionfalse
      \gdef\leq@current@section{#2}%
      \gdef\leq@current@subsection{}%
      \leq@original@section[#1]{#2}%
    }
    % Redefinir \subsection
    \renewcommand{\subsection}{%
      \@ifstar{\leq@subsection@star}{\@dblarg\leq@subsection@nostar}%
    }
    \def\leq@subsection@star#1{%
      \global\LeqStarSectiontrue
      \gdef\leq@current@subsection{#1}%
      \leq@original@subsection*{#1}%
      \global\LeqStarSectionfalse
    }
    \def\leq@subsection@nostar[#1]#2{%
      \global\LeqStarSectionfalse
      \gdef\leq@current@subsection{#2}%
      \leq@original@subsection[#1]{#2}%
    }

    % ----- Restauración de valores Globales de estilo ----------%
    % Flags
    \newif\ifInFootnote
    \InFootnotefalse
    \pretocmd{\@footnotetext}{\InFootnotetrue}{}{}
    \apptocmd{\@footnotetext}{\InFootnotefalse}{}{}
    % Profundidad de listas (soporta anidadas)
    \newcount\MyListDepth \MyListDepth=0
    \AtBeginEnvironment{la}{\advance\MyListDepth by 1}
    \AfterEndEnvironment{la}{\advance\MyListDepth by -1}
    \AtBeginEnvironment{lnum}{\advance\MyListDepth by 1}
    \AfterEndEnvironment{lnum}{\advance\MyListDepth by -1}
    \AtBeginEnvironment{itemize}{\advance\MyListDepth by 1}
    \AfterEndEnvironment{itemize}{\advance\MyListDepth by -1}
    \AtBeginEnvironment{enumerate}{\advance\MyListDepth by 1}
    \AfterEndEnvironment{enumerate}{\advance\MyListDepth by -1}
    % Restauración diferida: hazlo al INICIO del próximo párrafo real
    \newif\if@pendingRestore
    \@pendingRestorefalse
    \newcommand{\RequestRestore}{\global\@pendingRestoretrue}
    \everypar{%
      \if@pendingRestore
        \global\@pendingRestorefalse
        \ifInFootnote\else
          \ifnum\MyListDepth=0
            \setlength{\parskip}{\GlobalParskip}%
            \setstretch{\GlobalBaseLineStretch}%
          \fi
        \fi
      \fi
    }
    % Al final de bloques:
    \AfterEndEnvironment{equation}{\RequestRestore}
    \AfterEndEnvironment{equation*}{\RequestRestore}
    \AfterEndEnvironment{la}{\RequestRestore}
    \AfterEndEnvironment{lnum}{\RequestRestore}
    % Al final de macros:
    \apptocmd{\eqblock}{\RequestRestore}{}{}
    \apptocmd{\imagenfit}{\RequestRestore}{}{}
    
\makeatother

% ===================== ToC propio con 4 niveles (único bloque) =====================
\makeatletter

% --- Escritores: añadimos una línea al .stc en cada cabecera numerada ---
% Guardamos las versiones actuales para llamar al comportamiento normal
\let\MyTOC@orig@section\section
\let\MyTOC@orig@subsection\subsection
\let\MyTOC@orig@subsubsection\subsubsection
\let\MyTOC@orig@paragraph\paragraph

% SECTION
\renewcommand{\section}{%
  \@ifstar{\MyTOC@section@star}{\@dblarg\MyTOC@section@nostar}%
}
\def\MyTOC@section@star#1{%
  \MyTOC@orig@section*{#1}% sin entrada al .stc
}
\def\MyTOC@section@nostar[#1]#2{%
  \MyTOC@orig@section[{#1}]{#2}% comportamiento normal (escribe al .toc)
  % Escribimos también nuestra línea al .stc (texto con \numberline)
  \addcontentsline{stc}{section}{\protect\numberline{\thesection}#2}%
}

% SUBSECTION
\renewcommand{\subsection}{%
  \@ifstar{\MyTOC@subsection@star}{\@dblarg\MyTOC@subsection@nostar}%
}
\def\MyTOC@subsection@star#1{%
  \MyTOC@orig@subsection*{#1}%
}
\def\MyTOC@subsection@nostar[#1]#2{%
  \MyTOC@orig@subsection[{#1}]{#2}%
  \addcontentsline{stc}{subsection}{\protect\numberline{\thesubsection}#2}%
}

% SUBSUBSECTION
\renewcommand{\subsubsection}{%
  \@ifstar{\MyTOC@subsubsection@star}{\@dblarg\MyTOC@subsubsection@nostar}%
}
\def\MyTOC@subsubsection@star#1{%
  \MyTOC@orig@subsubsection*{#1}%
}
\def\MyTOC@subsubsection@nostar[#1]#2{%
  \MyTOC@orig@subsubsection[{#1}]{#2}%
  \addcontentsline{stc}{subsubsection}{\protect\numberline{\thesubsubsection}#2}%
}

% PARAGRAPH
\renewcommand{\paragraph}{%
  \@ifstar{\MyTOC@paragraph@star}{\@dblarg\MyTOC@paragraph@nostar}%
}
\def\MyTOC@paragraph@star#1{%
  \MyTOC@orig@paragraph*{#1}%
}
\def\MyTOC@paragraph@nostar[#1]#2{%
  \MyTOC@orig@paragraph[{#1}]{#2}%
  % Si no numeras los \paragraph, cambia \theparagraph por texto plano sin \numberline
  \addcontentsline{stc}{paragraph}{\protect\numberline{\theparagraph}#2}%
}

% --- Impresor del índice propio (lee el .stc) ---
\newcommand{\MyTOC}{%
  \par\bigskip
  {\centering\bfseries\Large Índice de contenido\par}%
  \medskip
  \begingroup
    \parindent=0pt\parskip=2pt
    \@starttoc{stc}%
  \endgroup
}

\makeatother
% ===================== fin ToC propio =====================


%%%%%%%%%%%%%%


% --- Datos del tema ---
\title{Tema 3.A.11 -- Teoría de la Producción (I)\\
\large Caracterización de la tecnología de la empresa a corto y largo plazo. El conjunto de posibilidades de producción. La función de producción. Rendimientos locales y globales a escala. Elasticidad de sustitución. Producción conjunta.}
\author{Marcelo Ortega}
\date{Febrero 2026}
\newcommand{\TemaCode}{3A11}

\oldtitle{Tema 3.A.9. Teoría de la Producción}
\changerationale{
    Se desarrolla el título del tema para garantizar un contenido mínimo común en las exposiciones de los opositores.
}

\usepackage{soul}\sethlcolor{yellow}% resaltado amarillo: \hl{texto} (Panel TCEE, ⌃H)
\begin{document}


\maketitle
\changebox

\newpage

% --- Página de resumen y modificaciones ---
\puntosclave{ %% Escribir puntos clave Apuntes CECO
    }
\vspace{1cm}

\modificaciones{
    
    }

\newpage
\MyTOC
\newpage

\mylistofequations
\clearpage

\MyListOfFigures
\clearpage


\pagenumbering{arabic}
\setcounter{page}{1}


\section{Introducción}\label{intro}


\subsection{Contextualización}\label{context}


\subsubsection{Relevancia}\label{importance}


\subsubsection{Historia}\label{history}



\subsubsection{Problemática}\label{problem} 




\subsection{Estructura de la exposición}\label{structure}




\clearpage
\section{Tecnología}
\puntosclave{
    La definición de tecnología y la representación del conjunto tecnológico. Morfología del conjunto.
    
    La Función de Producción, restricción técnica definidora de la Frontera de Posibilidades de Producción. Método de producción eficiente. Relación Marginal de Sustitución Técnica y Elasticidad de Sustitución.
    
    La Producción conjunta: el Proceso Multiproducto.
}
Para construir la función de producción seguiremos la metodología empleada en otras ramas de la microeconomía, distinguiendo tres etapas:
\begin{la}
    \item Definición del conjunto tecnológico;
    \item Obtención de la función de producción a partir de los axiomas de dicho conjunto; y, 
    \item Obtención y estudio de las isocuantas.
\end{la}
	
\subsection{Conjunto tecnológico}

\subsubsection{Definición: Tecnología y Conjunto Tecnológico}
La tecnología es la suma de técnicas, habilidades, métodos y procesos utilizados en la producción de bienes o servicios o en el logro de objetivos. Una primera forma de definir la tecnología, por lo tanto, es mediante los \textbf{vectores netput} (y). Cada uno de estos vectores hace referencia a un cierto proceso productivo:

$$y=(y_1,y_2,\dots,y_s)\in \mathbb R^s$$

Cuyos componentes: (i) negativos indican cantidades de factores adquiridos por la empresa; (ii) positivos representan cantidades obtenidas de productos; y, (iii) juntos conforman lo que se viene a llamar \textit{conjunto tecnológico}, que se refiere al conjunto de todos los planes de producción que resultan posibles en función de la tecnología y recursos de la empresa.
	
El conjunto tecnológico de factores-productos posibles representa por tanto, las combinaciones permitidas por la tecnología existente, ya sea de manera particular (sector o industria) o de manera general (economía).

\subsubsection{Axiomas: Propiedas}
Vamos a comenzar por establecer una serie de axiomas sobre el conjunto tecnológico. Con el fin de dar una buena imagen del conjunto tecnológico, dada su importancia, revisaremos dos definiciones: la de SEGURA y la de MAS-COLELL. Además, dividiremos el conjunto de axiomas en dos: (i) relativos al conjunto de vectores posibles; y, (ii) los relativos a las propiedades de la producción.

\paragraph{Conjunto vectorial: vectores \textit{neputs}}
Axiomas sobre combinaciones posibles en  el conjunto de vectores input-output (vectores netputs). MAS-COLELL y SEGURA coinciden en definir los siguientes axiomas:
\begin{lnum}
    \item \textbf{No vacío}. Debe haber alguna combinación obtenible pues de lo contrario no podría existir actividadproductiva alguna	Axiomas sobre combinaciones posibles en  el conjunto de vectores input-output (vectores netputs);
    \item \textbf{Cerrado}. Incluye su frontera;
    \item \textbf{Eliminación gratuita}. Si una tarea productiva es técnicamente posible, lo es también cualquier otra que implique iguales o mayores cantidades aplicadas de los inputs e iguales o menores cantidades de los outputs. (Tercer cuadrante pertenece al conjunto);
    \item \textbf{\textit{No-free lunch}}. No es posible producir cantidades positivas de output sin la utilización de ningún input;
    \item \textbf{Posibilidad de inacción}. La inacción es posible (el origen pertenece al conjunto);
    \item \textbf{Aditividad o libertad de entrada}. Si dos vectores son técnicamente posibles, su suma también lo es. También es consecuencia de que el conjunto agregado de producción (el conjunto que describe los planes de producción factibles para la economía como un conjunto) debe satisfacer la aditividad siempre que exista libertad de entrada;
    \item Posibilidad de inacción – La inacción es posible (el origen pertenece al conjunto)
\end{lnum}

Por último, MAS-COLELL añade un sexto axioma adicional. \textbf{Irreversibilidad}, es imposible revertir un vector de producción tecnológicamente posible para transformas una cantidad de output en la misma cantidad de input que fue usada para su generación.

De esta forma, las combinaciones posibles del conjunto vectorial de \textit{netputs} o conjunto de posibilidades de producción se puederepresentar:

\imagenfit{CPP_Axiomas.png}{Conjunto de Posibilidades de Producción o conjunto de combinaciones posibles de vectores \textit{netpus}. Diferencias SEGURA (izq.) y MAS-COLELL (dcha.)}{Elaboración propia}

\paragraph{Operaciones vectoriales: proceso productivo}
Por otro lado, debemos definir la axiomática relativa a las operaciones que se pueden realizar con los vectores \textit{netputs}. Es decir, una vez hemos definido el conjunto de posibilidades factibles de los vectores \textit{neputs} debemos definir como interaccionan/operan entre ellos. Es lo que se conoce como axiomática y producciones a escala.

\begin{lnum}
    \item \textbf{Aditividad o libertad de entrada}. Si dos vectores son técnicamente posibles, su suma también lo es. También es consecuencia de que el conjunto agregado de producción (el conjunto que describe los planes de producción factibles para la economía como un conjunto) debe satisfacer la aditividad siempre que exista libertad de entrada;
    \item \textbf{Non-increasing returns to scale} (Divisivilidad del proceso productivo). Dado un vector de netputs tecnológicamente posibles (perteneciente al conjunto tecnológico), sea cualquier número λ o α∈[0,1], entonces el vector escalado también pertenece al conjunto tecnológico. (Para MAS-COLELL et. al., este axioma es alternativo, puede o no presentarse, se combina por tanto con los siguientes dos axiomas particulares suyos);
    \item \textbf{Non-decreasing returns to scale} La tecnología exhibe esta propiedad si para cualquier escalar $\lambda\ge1$, un vector de netputs que pertenece al conjunto escalado forma otro vector de netputs que pertenece al conjunto tecnológico;
    \item \textbf{Constant-returns to scale}. Resulta de la únion de los Axiomas (2 y 3 del proceso productivo). Para cualquier escalar $\lambda\ge0$, un vector de netputs escalado formará otro vector perteneciente al conjunto;
    \item \textbf{Convexidad}. El conjunto de producción es convexo. Una combinación lineal de vectores netputs pertenece también al conjunto tecnológico, sea $\alpha\in[0,1]$. Se puede interpretar como que las combinaciones no equilibradas de inputs son más productivas que las combinaciones equilibradas;
    \item \textbf{Convexidad del conjunto}. El conjunto tecnológico es un cono si y so solo si satisface la propiedad de aditividad y la condición de non-increasing returns to scale. Se trata de la unión de los Axiomas (2 y 5 del proceso productivo) 
\end{lnum}
		
Como se puede entrever, existen incompatibilidades dentro de esta segunda relación aximomática. En concreto:
\begin{la}
    \item \textbf{Rendimientos constantes a escala}. Si hay constant-returns to scale implica que exigimos el A7 (aditividad) y A8 (divisibilidad), que implican directamente A5 (convexidad, pero no A6). Se puede apreciar en el Gráfico a:
    
    El punto $c$ se puede alcanzar mediante la transformación lineal del punto $a$ y $b$. El punto $c\in \bar Y$, y $c\not\in \text{Int}(Y)$. Por tanto, se incumple A6;
    \item \textbf{Rendimientos crecientes a escala}. Si hay increasing-returns to scale implica A7 (aditividad) y, en consecuencia, incumplimiento de A6 (convexidad, pero no tiene porque incumplir A5, isocuantas convexas). Se puede apreciar en el Gráfico b.:

    El punto $c$ se puede alcanzar mediante transformación lineal de $a$ y $b$, $\{a,b\}\in\bar Y$, pero $c\not\in\bar Y$ ni tampoco $c\not\in\text{Int}(Y)$. Por tanto, se incumple el A5, y también en consecuencia el A6, pero no el A1 ya que la combinación de dos planes pertenecerá a $\text{Int(Y)}$; y,
    \item \textbf{Rendimientos decrecientes a escala}. Si hay decreasing-returns to scale, entonces tenemos el A2 y A6 (y A5) pero se incumple directamente el A1. Se puede apreciar en el Gráfico c.:

    Dados dos planes/vectores de producción $\{a,b\}\in\bar Y$, si escalamos el vector con $\alpha\in[0,1]$, el vector resultante $c\in\text{Int}(Y)$ pero la combinación de $a$ y $b$ no lo estará, ya que sería el punto $k$ y este $k\not\in Y$. Por tanto, se incumple el A5.
\end{la}

\imagenfit{Axiomas_ProcesoProductivo.png}{Incompatibilidades y significado de la axiomática relativa al proceso productivo. Operaciones vectoriales dentro del conjunto tecnológico}{Elaboración propia}
 
\subsection{Función de Producción}
Los axiomas anteriores nos permiten representar la tecnología mediante una función de transformación o producción, que se define como una relación puramente técnica entre insumos productivos y volúmenes de producción.\footnote{
En realidad, podemos distinguir dos categorías de axiomas:
\begin{la}
    \item Necesarios (garantizan que la tecnología pueda representarse mediante una función de producción):
    \begin{lnum}
        \item No vacío;
        \item Cerrado;
        \item Eliminación gratuita; y,
        \item Imposibilidad de producción gratuita.
    \end{lnum}; y,
    \item Deseables (dan pie a una tecnología de buen comportamiento y, por ende, a una función neoclásica de buen comportamiento, por lo que son condiciones indispensables para la competencia perfecta):
    \begin{lnum}
        \item Posibilidad de inacción. Se puede elegir no producir nada; y,
        \item Convexidad. La convexidad implica la posibilidad de divisibilidad de escala (i.e. se puede reducir la escala de la producción). Además implica rendimientos a escala no crecientes. Los procesos extremos son menos productivos que los promediados. En el caso de un solo output este axioma implica que la función de producción es cóncava. Para maximizar el beneficio esta condición es esencial.
    \end{lnum}
\end{la}
} Recoge el máximo producto que se puede obtener dado un vector de factores productivos, corrresponde por tanto, en puridad, a la relación “output eficiente”, esto es, la Frontera de Posibilidad de Producción (frontera del Conjunto Tecnológico):

$$f(z_1,z_2,\dots,z_i)=(q_1,q_2,\dots,q_j)$$

Esta función de producción, cuyo objeto es maximizar el nivel de producción posible para cada vector de cantidades aplicadas de factores puede ser de:
\begin{la}
    \item Producción conjunta (varios inputs y varios outputs); o,
    \item Producción simple (varios inputs y un único output).
\end{la}

\subsubsection{La Función de Producción Neoclásica de Buen Comportamiento}
La función de producción simple (uniproducto) definida en el conjunto tecnológico que asocia un vector de inputs a una cantidad de un output fue propuesta inicialmente por PHILIP WICKSTEED en 1894:

\eqblock{
\begin{aligned}
    q=f(\vec{z})
\end{aligned}
}{Función de Producción uniproducto: Neoclásica de Buen Comportamiento. Tecnología: Función de Producción}

\subsubsection{Propiedades}
Con los axiomas mencionados hasta aquí, podemos garantizar que esta función de producción cumple las siguientes características:
\begin{lnum}
    \item Es de carácter cardinal a diferencia de la función de utilidad que obteníamos en el caso del problema del consumidor.\footnote{%
        Una función es cardinal si el valor númerico mismo tiene sentido o significado cuantitativo por lo que no solo la ordenación que hace de sus variables importa. Por eso cualquier tranformación monotona creciente NO vale para la función de producción, solo “valdrán” transformaciones lineales positivas (escala o traslación) puesto que estas si preservan la información.
    };
    \item Es no-decreciente, por el axioma de eliminación gratuita (las Productividades Marginales son no negativas); y,
    \item Además, será conveniente suponer que la función de producción es continua y dos veces diferenciable (esto es conveniente desde un punto de vista matemático para poder utilizar herramientas de cálculo diferencial).\footnote{%
        Este es un supuesto instrumental para el que es decisivo postular la continuidad de la sustituibilidad entre los factores y la infinitud de técnicas, así como que el bien que se produce sea perfectamente divisible.
    }
\end{lnum}
	
\subsubsection{Curvatura}
Sin embargo, todavía no podemos afirmar nada sobre su curvatura, que nos dará información sobre dos cosas:
\begin{lnum}
    \item La complementariedad de los factores productivos. Supondremos que existe cierto grado de complementariedad o sustituibilidad entre los factores productivos;\footnote{%
        De forma intuitiva, si suponemos dos factores productivos, trabajo y capital, son complementarios en cierto grado si solo podemos producir poco output si uno de los factores es escaso aunque el otro sea abundante. En este sentido, ambos factores productivos son importantes para la producción y la media de dos vectores de producción extremos (uno con capital elevado y trabajo reducido y otro al revés) producirá estrictamente más producto que al menos uno de los dos vectores iniciales extremos (JEHLE y RENY, 2011). ADAM SMITH ya hablaba de complementariedad: los trabajadores son más productivos cuanto más capital haya. Por ejemplo, en una función de producción de tipo Cobb-Douglas existe la idea de complementariedad (la PMg del capital aumenta cuanto aumenta el capital).
    } y,
    \item Las características de escala de la tecnología, es decir la curvatura global,\footnote{%
        Es muy importante diferenciar entre la curvatura global y la curvatura local de la función de producción.
    } que puede ser:
    \begin{la}
        \item Si el conjunto tecnológico fuese convexo, podremos afirmar que la función de producción será cóncava y, en el largo plazo, presenta rendimientos constantes a escala;
        \item Si el conjunto tecnológico fuese estríctamente convexo, podremos afirmar que la Función de Producción será estríctamente cóncava y, en el largo plazo, presentará rendimientos decrecientes a escala;\footnote{%
            En cualquier caso, no debemos olvidar que la función de producción se encuentra ligada con la función de costes siendo una parte del problema general de optimización a resolver y que desde una perspectiva conjunta de la teoría de la producción y costes nos haría tener en cuenta los precios y las limitaciones en las dotaciones de factores.
        }
        \item Si el conjunto tecnológico presenta rendimientos crecientes a escala, puede o no ser convexo en su frontera, entonces podremos afirmar que la función de producción es cuasi-cóncava y , en el largo plazo, presentará rendimientos crecientes a escala.
    \end{la}
\end{lnum}

\imagenfit{FPNBC_Tipos.png}{Tipos de Función de Producción neoclásicas de buen comportamiento. Rendimientos a escala}{Apuntes Victor Gutierrez. Tema 3.A.11}
 
Si se cumplen todas estas características (es decir, las primeras y solo alguna de escala), diremos que estamos ante una Función de Producción de Buen Comportamiento (luego levantaremos esta restricción para el análisis por FASES de CASSEL).

\subsubsection{La representación gráfica: Isocuantas}
Llamaremos \textbf{isocuantas} a los conjuntos de nivel de la función de producción. Esto es, las combinaciones entre factores que dan lugar a un mismo nivel de producto:

\imagenfit{FPNCB_Tipos_Isocuantas.png}{Proyección de los cortes (niveles) de la Función de Producción: Isocuantas}{Apuntes Víctor Gutierrez. Tema 3.A.11}
 
Analizaremos la Función de Producción ahora como función de insumos. Se trata por tanto de la propiedad local de sustitución de los insumos.  

\paragraph{Propiedades}
Dicho esto, las curvas isocuantas poseen las siguientes propiedades:
\begin{lnum}
    \item Forman un conjunto infinito pero numerable;
    \item No se cortan, por el axioma de eliminación gratuita;\footnote{%
        La monotonicidad de la función de producción (o eliminación gratuita / free-disposal) implica que: “Si usamos más insumos, nunca produciremos menos output”. En consecuencia, aunque pueda parecer que dependa de los rendimientos a escala que presente el conjunto tecnológico, pensemos que: si las isocuantas se cortaran, querría decir que un mismo vector de insumos produciría dos niveles de output diferentes. En todo caso, el incremento hará que parezca que se vayan acercando o alejando, pero nunca se cortaran.
    }
    \item Las curvas isocuantas son de clase $\mathcal C^n,\forall n>2$;
    \item Cuando más alejadas, más cantidad. Sin embargo, no será tan fácil como medir su distancia euclideana puesto que deberemos tener en cuenta el tipo de rendimientos a escala que presenta;
    \item Las curvas isocuantas son convexas (si asumimos que los factores de producción son cooperativos, complementarios);
    \item Proporcionan un ordenamiento cardinal de los distintos niveles de producción.
\end{lnum}

2.2.3.2. Relación Marginal de Sustitución Técnica (RMST)
Las curvas isocuantas tienen pendiente negativa, por el axioma de eliminación gratuita. Es decir, cuando se reduce la aportación de un factor, la única posibilidad de producir la misma cantidad del bien es usando más del otro factor.

La pendiente de las curvas isocuantas se obtiene diferenciando totalmente la función de producción y es lo que se conoce como Relación Marginal de Sustitución Técnica. En palabras, es la medida de la variación que debe experimentar  el uso de un $z_i$ ante la disminución de otro $z_j$ para mantener la producción a un nivel constante, por lo que hablamos, en cierto modo, de la medida del grado de sustituibilidad que existe entre los diferentes factores productivos, dado un nivel de producción determinado. Analíticamente, podemos definir la RMST de la siguiente forma:\footnote{%
    Habitualmente, la RMST se suele definir en valor absoluto, ya que normalmente consideraremos que la función de producción es no decreciente en los factores de producción debido al axioma de eliminación gratuita y por lo tanto la RMST es negativa (i.e. la pendiente de las curvas isocuantas es negativa).
}

\eqblock{
\begin{aligned}
    |{RMST}_L^K|=-\left.\frac{\mathrm d z_K}{\mathrm dz_L}\right|_{
    \begin{aligned}
        &\mathrm dq=0\\
        &\mathrm z_i=0;\;\;\forall i\neq K,L
    \end{aligned}}=\frac{\partial f/\partial z_L}{\partial f/\partial z_K}=\frac{{PMg}_{z_L}}{{PMg}_{z_K}}
\end{aligned}
}{Función de Producción: Relacción Marginal de Sustitución. Tecnología: Función de Producción}
 
Presenta la siguientes características:
\begin{lnum}
    \item Como la $RMST$ es la pendiente de la curva isocuanta en un punto, equivale al cociente de las productividades marginales ($PMg$); y,
    \item La concavidad (estricta o no)/ cuasi-concavidad de la función de insumos, esto es la convexidad de las isocuantas, implica que la $RMST$ disminuye en valor absoluto a medida que descendemos a lo largo de ésta.
\end{lnum}

\paragraph{La Elasticidad de Sustitución: ES}
Las unidades de medida de los diferentes factores productivos son diferentes, por tanto, la $RMST$ no podrá, en cierto modo, homogeneizarse. Para resolver este problema y poder dar con una forma de medir realmente el grado de sustituibilidad deberemos introducir la elasticidad de sustitución , que se define como el cociente entre la variación porcentual de la proporción entre dos factores (intesidad factorial) y la variación porcentual de la RMST:\footnote{%
    El concepto de elasticidad de sustitución lo introduce JOAN ROBINSON en el contexto de la controversia del capital de Cambridge [ver anexo A.2].

La Elasticidad de Sustitución para una tecnología dada se puede estimar mediante una regresión lineal de Mínimos Cuadrados Ordinarios de tal forma que:
$\text{factor-intesinty} = \beta_0 + \beta_1 · RSMT + \varepsilon$ ; donde,
tanto factor_intensity como rsmt se introducen en logaritmos y potencialmente podrían incluirse controles para solucionar problemas de endogeneidad y mejorar la precisión. b1 sería en este caso la estimación de la elasticidad de sustitución.
}

\eqblock{
\begin{aligned}
    \text{Elasticidad de Sustitución}=\sigma=\frac{\frac{\mathrm d(z_K/z_L)}{(z_K/z_L)}}{\frac{\mathrm d|{RMST}_L^K|}{|{RSMT}_L^K|})}=\frac{\mathrm{d}(z_K/z_L)}{\mathrm |{RMST}_L^K|}\cdot \frac{|{RMST}_L^K|}{(z_K/z_L)}
\end{aligned}
}{Elasticidad de sustitución: Relación Marginal de Sustitución e insumos productivos. Tecnología: Función de Producción}
  
De esta forma, conseguimos un número puro, independiente de las unidades de medición que ofrece una medida neutral del grado de sustituibilidad, esto es, de la curvatura de la isocuanta de un determinado nivel de output. 

Cuanto mayor sea esta expresión, más fácil es sustituir un factor por otro, en el sentido de que hace falta una menor variación de la RMST para acomodar un cambio en la proporción de los factores. 

Intuitivamente, una elevada elasticidad de sustitución supone una escasa sensibilidad de las PMg ante cambios en las proporciones de los factores (p.e. la elasticidad de sustitución de una función de producción con factores sustitutivos perfectos es igual a $+\infty$). Por el contrario, una baja elasticidad de sustitución es la respuesta a una importante sensibilidad de las PMg a cambios en las proporciones de los factores (p.e. la elasticidad de sustitución de una función de producción de coeficientes fijos o de Leontief es igual a 0).
Un tipo muy concreto de función que se usa habitualmente en la literatura es la función de producción de Elasticidad de Sustitución Constante (Constant Elasticity of Substitution, CES) . Se refiere a un tipo particular de función que combina dos o más tipos de factores de producción en una cantidad agregada. Esta función agregada exhibe una elasticidad de sustitución constante.\footnote{%
    La función de producción CES fue introducida por ROBERT SOLOW en su popular obra “A contribution to the theory of economic growth” (1956) [ver Tema 3.A.43] y posteriormente popularizada por ARROW, CHENERY, MINHAS y SOLOW (1961).
}

\eqblock{
\begin{aligned}
    q=A\cdot\left[\sum_{i=1}^ns_i\cdot z_i^{\frac{\sigma-1}{\sigma}}\right]^{\frac{\sigma}{\sigma-1}}
\end{aligned}
}{Elasticidad de sustitución: Elasticidad de sustitución constante. Tecnología: Función de producción}
 
Dentro de este tipo de funciones de producción se encuentran algunas de las más utilizadas por la literatura (aquí especificadas para el caso de dos factores productivos), entre otras:
\begin{lnum}
    \item Si la elasticidad de sustitución es igual a cero: se trata de una función de producción de coeficientes fijos (o de complementarios perfectos o de tipo Leontief):
    $$q=\min{\left\{\dfrac{z_K}{a}; \dfrac{z_L}{b}\right\}}$$
    \item Si elasticidad de sustitución es igual a uno: converge a una función de tipo Cobb-Douglas:
    $$q=A\cdot z_K^\alpha \cdot z_L^\beta;\quad \text{donde, }\{A,\alpha, \beta\}\in\mathbb R^+$$
    \item Si elasticidad de sustitución tiende a infinito: converge a una función de producción linear (o de sustitutivos perfectos):
    $$q=a\cdot z_K+b\cdot z_L$$
\end{lnum}
 
\subsubsection{Funciones de producción típicas}
Estas formas funcionales son muy utilizadas en la literatura económica, por ejemplo en ámbitos como el crecimiento económico.

\imagenfit{FP_CES.png}{Función de Producción de tipo CES}{Apuntes Víctor Gutierrez. Tema 3.A.11}

Por eso conviene verlas con más detalle.

\paragraph{Función COBB-DOUGLAS}
La función COBB-DOUGLAS es la forma más común en toda la literatura micro y macroeconómica. Aunque todas las funciones con forma COBB-DOUGLAS se caracterizan por mostrar elasticidades de sustitución constantes e iguales a uno, varían en los rendimientos a escala que muestran y en la contribución a la producción total de los diferentes factores de producción. 

Una característica fundamental de esta tecnología es que la PMg de los factores es siempre positiva: como resultado de ello, todo el plano es parte de la región económica de producción (no existen lineas de contorno), ya que las isocuantas son asintóticas a los ejes. Esta función pertenece, por tanto, a la categoría de funciones de “buen comportamiento” y su forma más común es:

\eqblock{
\begin{aligned}
    &\text{Forma general:}&&Y=A\cdot\Pi_{i=1}^nz_i^{\alpha_i}\qquad \text{donde,} \{A, \alpha_i\}\in\mathbb R^+\\[1em]
    &\text{Forma convencional:}&&\boxed{Y=A\cdot K^\alpha\cdot L^{1-\alpha}}\;\Longrightarrow\;\underset{\substack{\text{\scriptsize Elasticidad de}\\\text{\scriptsize Sustitución constante}}}{\sigma=1}
\end{aligned}
}{Formas comunes de la Función de producción: COBB-DOUGLAS}

Gráficamente, podríamos representar sus isocuantas y los rendimientos marginales de un insumo productivo ($L$) de la siguiente forma (teniendo en cuenta únicamente dos factores productivos):

\imagenfit{COBB-DOUGLAS_FP.png}{Isocuantas de la Función COBB-DOUGLAS y rendimiento marginal del factor productivo trabajo ($L$)}{Elaboración propia}
 
\paragraph{Función de complementaridad perfecta: Tecnología LEONTIEF}
En este tipo de funciones los inputs contribuyen a aumentar la producción en una proporción determinada y constante. Las funciones de Leontief son la pieza fundamental de los modelos Input-Output y, permite modelar contextos en los que resulta razonable asumir rigidez importante de los procesos productivos. Su forma funcional es:

\eqblock{
\begin{aligned}
    &\text{Forma general:}&&Y=\min\left\{\frac{z_1}{a_1}, \frac{z_2}{a_2}, \cdots, \frac{z_n}{a_n}\right\}\qquad \text{donde,}\;\; a_i>0\\[1em]
    &\text{Forma convencional:}&&\boxed{Y=\min\left\{\frac{z_K}{a_1}, \frac{z_L}{a_2}\right\}}
\end{aligned}
}{Formas comunes de la Función de producción: Coeficientes fijos (LEONTIEF)}

Podemos representar gráficamente las isocuantas y las relaciones de complementariedad:

\imagenfit{FP_LEONTIEFF.png}{Relación de complementariedad e isocuantas quebradas. Función de Producción tipo LEONTIEFF}{Elaboración propia}
 
Otra manera de recoger rigideces en la sustitución entre actores es a través de la combinación de varías funciones LEONTIEF, lo que da lugar a isocuantas quebradas (dcha.), para producir un bien existen sólo unos pocos procesos de producción. En este caso existen tres procesos de producción para obtener el mismo nivel de producción, con distintas proporciones de factores. Se puede considerar que se trata de tres procesos de producción en proporciones fijas (P1, P2 y P3). Son las isocuantas que más se acercan a la realidad del corto plazo y son utilizadas en la programación lineal.

\paragraph{La función de elasticidad de sustitución constante (CES)}
La función CES es una generalización de las anteriores: representan casos particulares de CES con un valor determinado para la elasticidad de sustitución. La función CES se presente en su forma más habitual: la forma homotética.

De manera general, se suele representar con isocuantas continuas y convexas supone la sustituibilidad continua y se supone que el número de procesos de producción aumenta indefinidamente. La cantidad del bien producido se mide por la distancia al origen de coordenadas a algún punto de la isocuanta. Son las más empleadas en la modelización micro y macroeconómica, pos sus evidentes ventajas analíticas. Su forma general es:

\eqblock{
\begin{aligned}
    &\text{Forma general:}&&Y=A\cdot \left[\sum_{i=1}^n\delta_i\cdot z_i^{-\rho}\right]^{-\frac{v}{\rho}}\qquad \text{donde,}\;\;\{v,\delta_i,A\}\in \mathbb R^+\quad \text{y}\quad \sum_{i=1}^n\delta_i=1\;\;\wedge\;\;\rho\ge-1\\[1em]
    &\text{Forma convencional:}&&\boxed{Y=\left(K^{\frac{1}{2}}+L^{\frac{1}{2}}\right)^2}\;\Longrightarrow\;\underset{\substack{\text{\scriptsize Elasticidad de}\\\text{\scriptsize Sustitución constante}}}{\sigma=2}
\end{aligned}
}{Formas comunes de la Función de producción: CES}

\subsubsection{Extensiones}
Hemos visto que existen dos notaciones para representar la restricción tecnológica de la empresa:
\begin{la}
    \item Explícita: $y \leq f(z_1, z_2)$; y,
    \item Implícita: $y-f(z_1, z_2) = g(z_1, z_2, y) \leq 0$
\end{la}

Si se produce de manera eficiente, ambas se satisfarían con igualdad (Frontera de Posibilidades de Producción). Trataremos, brevemente la segunda de las notaciones para explicar la empresa multiproducto y emplearemos el vector $y$ como netputs. Además, asumiremos que: (i) la producción de los bienes es complementaria; y, (ii) la tecnología permite la sustitución de los productos.

\paragraph{Producción Multiproducto: Curva de Transformación}
Esta tecnología se define como aquélla en la que el número de netputs de signo positivo (outputs) en su método de producción es mayor que uno (esto es, existen diferentes outputs). 

Es decir, la introducción de la empresa multiproducto no supone ningún cambio sustancial al análisis ya realizado de las tecnologías de producción de la empresa. Como hasta ahora, la FPP es el objeto más relevante para modelizar la restricción técnica de la empresa, y su morfología o pendiente contiene información relevante. Podemos recoger el significado económico de la pendiente en un método de producción consistente en un vector de dos neputs como sigue. 

Supongamos la función de producción $g$ (en forma explícita) que viene caracterizada por:

$$g(z_n, y_k)\leq0\qquad\text{donde,}\;\;\forall n,k\in\mathbb Z$$

Sean $z_i,z_j$ insumos del proceso de producción. Tenemos tres definiciones esenciales en este escenario:

\eqblock{
\begin{aligned}
    &\underset{\text{\scriptsize Representa el CPP del proceso productivo global, espacio vectorial de posibilidades}}{\text{Relación Marginal de Sustitución Ténica:}}&&\frac{\mathrm dz_i}{\mathrm dz_j}=\frac{(\partial g/\partial z_i)}{(\partial g/\partial z_j)}=-{RMST}_i^j\\[1em]
    &\text{Productividad Marginal del insumo $z_i$ en $y_i$:}&&\frac{\mathrm dy_i}{\mathrm dz_i}=\frac{(\partial g/\partial y_i)}{(\partial g/\partial z_i)}=-{PMg}_{z_i}\\[1em]
    &\underset{\text{\scriptsize Represente la Curva de Transformcación entre outputs}}{\text{Relación Marginal de Transformación:}}&&\frac{\mathrm dy_i}{\mathrm dy_j}=\frac{(\partial g/\partial y_i)}{(\partial g/\partial y_j)}=-{RMT}_i^j\\[1em]
\end{aligned}
}{Función de Producción Multiproducto: Relaciones. Tecnología: Función de Producción}

Gráficamente, estas tres relaciones se representarían de la siguiente forma:

\textbf{INCLUIR GRÁFICOS RELACIONES}

Es facil comprobar, en el caso multiproducto, que los incrementos en un input requieren  reducciones en otro input. Alternativamente, incrementos en un input manteniendo el resto constante requieren de incrementos en el uso de consumo de inputs en el proceso productivo. Es decir, cada curva de transformación (c) corresponde a un nivel de consumo de inputs determinado y el incremento en el nivel de inputs consumidor desplaza la curva de trasformación desde el origen hacia fuera.\footnote{%
    Si estuviésemos ante un vector de netputs con tres o más outputs también podría desplazarse la curva de transformación hacia fuera reduciendo consecuentemente la producción del otro/s bien/es.
} En otras palabras, las disminuciones en otros productos o los incrementos en los insumos permiten producir más tanto de un bien como de otro. 

\paragraph{Producción Conjunta}


\paragraph{Tecnologías aditivias y Proporciones Fijas}
Obtendríamos variaciones a esta curva de transformación cóncava al origen levantando los supuestos de partida, en concreto:
\begin{la}
    \item Levantando el (ii) supuesto -esto es, la tecnología no permite la sustitución perfecta entre productos-, las curvas de transformación serían rectangulares, lo que indica que un incremento en la producción de un bien requiere un incremento en el nivel de insumos y no puede lograrse reduciendo la producción del otro bien;\footnote{%
        Esto sucede, por ejemplo, en la industria química, donde los productos deben ser producidos en proporciones fijas de inputs.
    } y,
    \item Levantando el (i) supuesto -esto es, la complementariedad de los dos outputs-, nuestra curva de transformación sería una recta. En estos casos, en realidad estamos ante tecnologías aditivamente separables o independientes, ya que se puede modelizar la producción de cada output con su propia función de producción. 
\end{la}

Habiendo estudiado la función de producción y habiendo estudiado algunos ejemplos de funciones de producción vamos a pasar al segundo bloque de la exposición en el cual introduciremos herramientas analiticas para estudiar la función de producción, comentando las leyes de la función de producción en el corto, largo y muy largo plazo.


\clearpage
\section{Leyes de producción}
El análisis por horizontes temporales muestra la dimensión intertemporal de la Teoría de la Producción. Fue desarrollado por ALFRED MARSHALL en su «Principles of Economics» (1890), y se ha mantenido prácticamente inalterado desde entonces. Distinguiendo entre diferentes horizontes temporales, podemos introducir otros conceptos fundamentales en la teoría de la producción:
\begin{la}
    \item \textbf{Corto plazo}. Nos encontramos ante un escenario en el que existe rigidez en la oferta de algún factor de producción (principalmente, el capital), pudiendose considerar fijo. Otros factores de producción (por ejemplo, el trabajo) presentan una oferta menos rígida en el corto plazo. Para el análisis del corto plazo, buscaremos los conjuntos de nivel de aquél factor rígido, en concreto, querremos analizar el comportamiento del resto de factores (su PMg) y como varían el output, dado un nivel de dicho factor;
    \item \textbf{Largo plazo}. Nos encontramos ante un escenario en el que todos los factores son variables, la oferta es variable. En este caso, centraremos nuestro análisis y planeamiento sobre el mapa de isocuantas, enfocándonos en ver cómo reacciona la producción ante cambios equiproporcionales en ambos factores productivos. En otras palabras, estudiaremos los rendimientos a escala; y,
    \item \textbf{Muy largo plazo}. Nos encontramos ante un escenario en el que, además de una oferta elástica de factores de producción, la tecnología es también variable. En términos del análisis gráfico, en el muy largo plazo, el conjunto tecnológico varía y la función de producción se desplaza hacia arriba. Este análisis nos permite introducir el concepto de progreso tecnológico.
\end{la}
	
Presentamos una división temporal del proceso de producción, veremos que en cada marco temporal, la importancia reside en el análisis de los factores productivos variables, cómo encontrar la combinación óptima para cada momento temporal y diferentes alternativas para comprender el progreso tecnológico a largo plazo.

\textbf{NOTA.-} \textit{Introduciremos y hablaremos sobre la región económica y el Análisis por FASES de CASSEL, por tanto, es imprescindible comentar que , y se incumplirá el axioma de eliminación gratuita. Esto es así porque: suponemos que a partir de un punto al aumentar el factor de producción variable el nivel de producción disminuye, es decir, el problema no está en la PMg decreciente , sino en la productividad marginal negativa a partir de cierto nivel. Relajaremos este supuesto ya que estudiaremos el análisis por etapas desarrollado por GUSTAV CASSEL en 1918 , mientras que la axiomática fue desarrollada con posterioridad.} REFORMULAR

\imagenfit{Esquema_TiempoProduccion.png}{¿Qué diferencias hay entre cada periodo? Esquema de variaciones y leyes que rigen el proceso productivo atendiendo a los factores productivos y las propiedades de la Función de Producción}{Apuntes Víctor Gutierrez. Tema 3.A.11}

\subsection{Corto plazo: Factor fijo}
Empezando por el análisis a corto plazo, es necesario advertir que vamos a utilizar una función de producción que va a incumplir uno de los axiomas que hemos mencionado antes, en concreto, va a incumplir el axioma de eliminación gratuita, ya que la función de producción no va a ser monótona y creciente y las curvas isocuantas no van a ser decrecientes en todos sus tramos.\footnote{
Las productividades marginales podrían ser decrecientes pero que su variación tendiera a un número no negativo. Esto sucede con las funciones de producción que cumplen las condiciones de Inada y que son muy habituales en la literatura:
\begin{lnum}
    \item Función de producción continua y diferenciable;
    \item Función de producción estríctamente creciente en factores;
    \item El límite de la función cuando cualquier factor tiende a cero es cero (los Factores no son indendientes entre sí);
    \item El límite de la función cuando cualquier factor tiende a infinito es infinito;
    \item El límite a infinito respecto a cualquiera de los factores, de la derivada de la función respecto a cualquier factor, es también infinito; y,
    \item El límite a infinito respecto a cualquiera de los factores, de la derivada de la función respecto al otro factor es igual a cero. 
\end{lnum}

Algunos ejemplos de estas funciones son las funciones de tipo Cobb-Douglas. Con este tipo de funciones de producción la productividad marginal nunca sería negativa y toda la función formaría parte de la región económica de producción (no habría líneas de contorno, pues la RMST no alcanzaría nunca alcanzaría ni el valor cero ni el infinito).

Según cómo sea la función de producción, los rendimientos marginales decrecientes pueden aparecer a partir de determinada utilización de los factores en el punto de inflexión A o desde el principio u origen de coordenadas. En la exposición nos centraremos en el primer caso debido a su mayor generalidad.
 
\imagenfit{InflexionA_FP_CASSEL.png}{Productividad marginal del factor trabajo, manteniendo capital constante}{Apuntes Victor Gutierrez. Tema 3.A.11}
}

Por lo tanto, estamos ante una función de producción que no es de buen comportamiento:

\imagenfit{Grafico_CP_FP_CASSEL.png}{Obtención del gráfico para el análisis de corto plazo de la productividad marginal en el caso de una función de producción de rendimientos mixtos}{Apuntes Víctor Gutierrez. Tema 3.A.11}
 
La forma o curvatura de la función de producción a corto plazo se debe a los rendimientos mixtos que aparecen al mantener fijo un factor e ir aumentando el factor variable:
\begin{lnum}
    \item En un primer momento, a medida que aumenta el factor variable, la producción aumenta a un ritmo cada vez mayor;
    \item A partir de cierto punto, la producción crece a un ritmo cada vez menor  hasta que alcanza un máximo;\footnote{%
        Se denomina Ley de los Rendimientos Marginales Físicos Decrecientes porque la producción se mide en unidades físicas, o también se denomina Ley de las Proporciones Variables porque la proporción que guardan entre sí el factor fijo y el variable va cambiando a lo largo de toda la función de producción.
    } y,
    \item A partir de ese máximo, si sigue aumentando el factor variable, la producción disminuye.
\end{lnum}

¿Y por qué utilizamos esta función de producción si no es habitual en la práctica e incumple los axiomas que hemos impuesto sobre la tecnología? 

\subsubsection{Productividad}
Vemos como con esta función de producción, si mantenemos constante el nivel de capital, a medida que aumenta el factor variable, la producción aumenta a un ritmo cada vez mayor. Pero a partir de cierto punto la producción crece a un ritmo cada vez menor hasta que alcanza un máximo. Y a partir de ese máximo, si sigue aumentando el factor variable, el incremento del producto se torna negativo. Esto se puede medir en términos de Productividades, es decir, el rendimiento que proporcionan las variaciones del factor variable (elástico) sobre el output del proceso. Introduciremos la Productividad Marginal (PMg) y la Productividad Media (PMe), teniendo en cuenta que hay relación entre estas.\footnote{%
    La PMe y la PMg tienen el mismo valor cuando la PMe es máxima. Esto se puede apreciar claramente de forma analítica, ya que, por la Ley de los rendimientos decrecientes, veremos como, si derivamos la PMe, su máximo se obtendrá cuando la derivada de PMe sea igual a cero. Por otro lado, la PMg se mide como la derivada de la función de producción sobre el factor que se evalúe, por lo tanto, en el  máximo, la PMe será igual a la PMg.
}

\paragraph{Productividad Media: PMe}
La productividad media de un factor de producción variable (PMe) es igual al producto total dividido por la cantidad de factor variable utilizado y expresa la relación existente entre el nivel de producción obtenido y la cantidad de factor variable aplicado:

\eqblock{
\begin{aligned}
    PMe_L=\frac{Q}{L}
\end{aligned}
}{Productividad media de un factor productivo: Corto plazo. Leyes de producción: Análisis por fases de CASSEL}
 
\paragraph{Productividad Marginal: PMg}
La productividad marginal (PMg) es el incremento que se produce en el producto total al aumentar en una unidad el factor variable, permaneciendo constante el factor fijo. Es por ello, la derivada de la función de producción en cada uno de sus puntos. Se refiere siempre a la última unidad de factor variable incorporado al proceso productivo:

\eqblock{
\begin{aligned}
    PMg_L=\frac{\mathrm dQ}{\mathrm dL}
\end{aligned}
}{Productividad marginal de un factor productivo: Corto plazo. Leyes de producción: Análisis por fases de CASSEL}

\subsubsection{Etapas de la producción}
La definición de todos estos conceptos previos nos permite ahora introducir uno de los resultados más importantes, y bellos, del análisis sobre el corto-plazo. Siguiendo a CASSEL (1918), podremos identificar tres etapas en esta función de producción:
\begin{lnum}
    \item \textbf{FASE I}. Donde la productividad media y marginal aumentan con cada unidad adicional, que en el primer diagrama corresponde al tramo hasta el punto de inflexión ($B$). La fase I termina en el margen extensivo (u óptimo técnico), donde se alcanza el máximo del Producto Medio. En esta primera etapa, el PMe es creciente y, además, la $PMg$ del factor de producción fijo es negativa, y viceversa. Existe demasiada cantidad de factor fijo en relación con el factor variable utilizado;
    \item \textbf{FASE II}. En la segunda etapa, que abarca desde el máximo del producto medio hasta el máximo técnico (donde la PMg del factor elástico es cero porque no se puede ya crecer de manera intensiva), es donde se debe producir (nótese que en este tramo es cuando nos encontramos dentro de la región económica de producción, sin situarnos en la linea de contorno). En esta etapa se da el equilibrio de la producción y las cantidades utilizadas de factores son las óptimas. La relación de factores de producción es eficiente en toda esta etapa; y,
    \item \textbf{FASE III}. Cuando la $PMg$ del factor elástico es negativa, entramos en la Fase III, donde se está sobreconsumiendo el factor variable (ya no hay posibilidad de aumentar la producción de manera intensiva).
\end{lnum}

\subsubsection{Elasticidad-producto del factor}
Todo esto está relacionado con la elasticidad-producto del factor, que mide la variación porcentual en la cantidad producida como resultado de una variación porcentual en el nivel de un factor de producción:

\eqblock{
\begin{aligned}
    \varepsilon_{Q,L}=\frac{(\mathrm dQ/Q)}{(\mathrm dL/L)}=\underset{PMg_L}{\frac{\mathrm{d}Q}{\mathrm dL}}\cdot \underset{1/PMe_L}{\frac{L}{Q}}=\frac{PMg_L}{PMe_L}
\end{aligned}
}{Elasticidad factor producto: Corto plazo. Leyes de producción: Análisis por fases de CASSEL}
 
Para el caso de nuestro ejemplo, podemos afirmar que:
\begin{lnum}
    \item En la FASE I: (1, inf.]:
    \item Margen Extensivo (óptimo técnico): {1};
    \item En la FASE II: (0,1);
    \item Margen Intensivo (máximo técnico): {0}; y,
    \item En la FASE III: (-inf, 0).
\end{lnum}

\subsubsection{Región Económica: Isoclinas}
Precisamente, estas “FASES” de la producción / productividad pueden ser representados sobre el plano cartesiano de factores de tal forma que:

\imagenfit{RegionEconomica_CP_CASSEL.png}{Región económica de producción}{Apuntes Víctor Gutiérrez. Tema 3.A.11; \href{policonomics.com/es/lp-produccion1-region-economica-produccion/}{Policonomics (Región económica de producción)}}

Dado el mapa de isocuantas, si mantenemos constante el factor inelástico (a un nivel fijo) y variamos el factor elástico, la producción va aumentando hasta el punto A, es decir, la productividad marginal de A en ese tramo es positiva. Pero si a partir de A se sigue aumentando el factor elástico, pasamos a curvas isocuantas asociadas a un nivel de producción cada vez menor, o sea, la PMg del factor elástico es negativa a partir de A.\footnote{%
    De la misma manera podríamos haber invertido las elasticidades y suponer que el factor elástico fuera el capital. En el cuerpo del tema hemos supuesto que el factor elástico es el trabajo y el factor inelástico es el capital ya que es el supuesto más habitual en la literatura económica.
}

El área en rojo representa la región económica, que muestra las diferentes combinaciones de factores de producción con un determinado coste, que tienen sentido económico, esto es, que se encuentran dentro de la FASE II de CASSEL, y por tanto, son técnicamente eficientes. Es decir, las zonas que están fuera de la región implican que al menos uno de los factores tiene PMg negativa (FASE I o FASE III), por lo que se podría producir más con una cantidad menor del bien con PMg negativa (en especial, será aquel factor que sea elástico, puesto que lo que nos interesará será optimizar la producción variando la cantidad del factor elástico). Esta región viene delimitada por las lineas de contorno (isoclinas), que son simplemente los límites más allá de los cuales uno de los factores está siendo usado en exceso. Por lo tanto:
\begin{la}
    \item Fuera de la región económica de producción, existe una clara ineficiencia, y la empresa estaría mejor usando menos de uno de los dos factores, lo que disminuiría los costes y aumentaría la producción; y,
    \item Por el contrario, dentro de la región económica, para mantener constante el nivel de producción, si varía la cantidad de un factor debe compensarse negativamente con la del otro. De lo contrario, si lo mantenemos constante y aumentamos la dotación de un factor, pasaríamos a otra isocuanta con un nivel de producción mayor.
\end{la}

Además, del gráfico también podemos ver donde se encuentran los puntos donde cada factor presente PMg = 0. Precisamente son los puntos que delimitan la región económica: (i) en la linea de controno superior, la $PMg_K = 0$; y, (ii) en la lineadecontrno inferior, la $PMg_L = 0$. 

Esto es, sobre la linea de contorno superior, el factor inelástico presenta PMg igual a cero y, sobre la linea de contorno inferior, el factor elástico presenta PMg igual a cero. Además, sabemos por la configuración de la función de producción que, sobre las lineas de contorno, la Relación Marginal de Sustitución Técnica (cociente de PMg’s) calculada sobre el Trabajo respecto al Capital, i.e. $PMg_L/PMg_K$:\footnote{%
    El lugar geométrico a lo largo del cual la RMST es constate se denomina isoclina y puede presentar cualquier forma según sea la función de producción. Si se trata de una función homotética, la isoclina es una línea recta ya que los factores de producción se modifican en la misma proporción, y la razón que guardan entre sí es constante (al igual que pasaba en el caso de la curva renta-consumo para el caso de la teoría de la demanda del consumidor [ver Tema 3.A.8]). Nótese que las líneas de contorno que definen la región económica de la producción son isoclinas, ya que a lo largo de estas la RMST es constante e igual a cero o infinito:

\imagenfit{RMST_CP_CP.png}{}{Apuntes Victor Gutierres. Tema 3.A.11}
}
\begin{la}
    \item \textbf{Cero}. Si estamos en la linea de contorno del factor elástico y (óptimo técnico); y,
    \item \textbf{Infinito}. Si estamos sobre la linea de contorno del factor inelástico (máximo técnico).
\end{la}
	
\subsection{Largo plazo: Factores variables}
Habiendo visto el análisis de corto plazo, pasamos al análisis del largo plazo, en el que todos los factores productivos son elásticos, es decir, un horizonte de planificación de la empresa, no uno de producción. 

Por tanto, en el largo plazo lo que se analizará será, valiéndose de las curvas isocuantas, la capacidad de sustituir factores de producción entre ellos.\footnote{Siempre sujeto, por supuesto, a la morfología concreta de cada tecnología. 
}

Nos centraremos en el concepto de rendimiento que aparece a la hora de cuantificar el impacto sobre el nivel de producción si varían todos los factores en la misma proporción.\footnote{%
    Este análisis ya lo hicimos cuando definíamos el conjunto de posibilidades de producción y la tecnología de producción. Aquí haremos hincapié en su representación y cuantificación.
} A esta variación equiproporcional en todos los factores de producción se le puede denominar variación en la escala de producción. 

Graficamente, los rendimientos a escala serían el plano bisector:

\imagenfit{Rendimientos_Ilustracion_R3.png}{Ilustración gráfica de diferentes tipos de rendimientos a escala en $\mathbb R^3$}{Apuntes Víctor Gutiérrez. Tema 3.A.11}

\subsubsection{Elasticidad-escala  o coeficiente de la función}
Matemáticamente, esto da lugar a la elasticidad-escala (o coeficiente de la función), que se define como el cociente entre la variación porcentual en la producción y la variación porcentual en las cantidades de todos los factores. Si tomamos todos los factores de producción como factores elásticos o variables, entonces:\footnote{%
    Como en el caso de la ES, una estimación empírica de esta elasticidad-escala es sencilla empleando una regresión lineal con logaritmos.
}

\eqblock{
\begin{aligned}
    &\mathrm dL/L=\mathrm dK/K=\mathrm dλ/λ,\qquad \text{siendo $\lambda$ la escala,entonces:}\\[1em]
    &\hspace{7em}\boxed{\varepsilon_{Q,\lambda}=\frac{\mathrm Q/Q}{\mathrm d \lambda/\lambda}=\frac{\mathrm dQ}{\mathrm \lambda}\cdot \frac{\lambda}{Q}}
\end{aligned}
}{Elasticidad a escala: Largo plazo. Leyes de producción: Análisis por fases de CASSEL}

Podremos, con ayuda de la elasticidad-escala de la función, clasificar los procesos productivos de la siguiente manera:
\begin{la}
    \item Diremos que el prcoeso de producción presenta rendimientos decrecientes a escala si su elasticidad-escala es inferior a 1;
    \item Diremos que el proceso de producción presente rendimientos constantes a escala si su elasticidad-escala es igual a 1; y,
    \item Diremos que el proceso de producción presenta rendimientos crecientes a escala si su elasticidad-escala es mayor a 1.
\end{la}

Hay que destacar que esta elasticidad no es un parámetro fijo (excepto en algunas formas funcionales, en especial la CES), sino que puede variar en distintos niveles de producción. Es decir, una función puede tener rendimientos crecientes, decrecientes y constantes en distintos subconjuntos de su tecnología. En el caso de que estos rendimientos varíen en el rango de los niveles de los factores, hablaremos de rendimientos locales a escala. En el caso de que estos sean invariantes, hablaremos de rendimientos globales a escala.

\paragraph{Relación entre productividades y rendimientos a escala}
Importante, pues ha sido objeto de mucha confusión.\footnote{%
    Por ejemplo, el uso de la \textit{Ley de rendimientos decrecientes}, tradicionalmente utilizada para referirse a los efectos en un factor variable y que no guarda nada que ver con los rendimientos a escala.
} Gráficamente, la productividad se mide comparando las isocuantas que se van alcanzando a medida que varía uno de los factores. Por su parte, los rendimientos a escala se miden comparando las isocuantas que se van alcanzando a medida que varían todos los factores proporcionalmente a lo largo de un radio que parte del origen. Analíticamente podemos demostrar la relación entre productividad y rendimientos, que puede obtenerse demostrando que la elasticidad-escala o coeficiente de la función es igual a la suma de las elasticidades-producto de los factores.\footnote{%
    Partimos de la siguiente función de producción: $Q=Q(L,K)$.

Si diferenciamos obtenemos que: $\mathrm dQ=\frac{\partial Q}{\partial L}\cdot\mathrm dL+\frac{\partial Q}{\partial K}\cdot \mathrm K=L\cdot\frac{\partial Q}{\partial L}\cdot\frac{\partial L}{L}+K\cdot\frac{\partial Q}{\partial K}\cdot\frac{\partial K}{K}$

Si dividimos por el outpur, obtenemos: $\frac{\partial Q}{Q}=\frac{L}{Q}\frac{\partial Q}{\partial L}\frac{\mathrm dL}{L}+\frac{K}{Q}\frac{\partial Q}{\partial K}\frac{\mathrm d K}{K}$

Como la variación que efectuamos es equiproporcional (lo mismo en cada factor), obtenemos: $\frac{\partial Q}{Q}=\frac{L}{Q}\frac{\partial Q}{\partial L}\frac{\mathrm d L}{L}+\frac{K}{Q}\frac{\partial Q}{\partial K}\frac{\mathrm d K}{K}$, donde $\frac{\mathrm dL}{L}=\frac{\mathrm K}{K}=\frac{\mathrm d\lambda}{\lambda}$. Entonces, obtenemos que: $\frac{\mathrm Q}{Q}=\frac{\mathrm d\lambda}{\lambda}\left[\frac{\partial Q}{\partial L}\frac{L}{Q}+\frac{\partial Q}{\partial K}\frac{K}{Q}\right]$.

Sabiendo que, $\varepsilon_{Q,L}=\frac{\partial Q/Q}{\mathrm d L/L}$ y $\varepsilon_{Q,K}=\frac{\partial Q/Q}{\mathrm d K/K}$. Tenemos que: $\frac{\mathrm d Q}{Q}=(\varepsilon_{Q,L}+\varepsilon_{Q,K})\frac{\mathrm d\lambda}{\lambda}\Longrightarrow\frac{\mathrm d Q/Q}{\mathrm d \lambda/\lambda}=\varepsilon_{Q,\lambda}=\varepsilon_{Q,L}+\varepsilon_{Q,K}$
}

\eqblock{
\begin{aligned}
   \varepsilon_{Q,\lambda}=\varepsilon_L+\varepsilon_K 
\end{aligned}
}{Elasticidad a escala: Descomposición. Leyes de producción: Análisis por fases de CASSEL}

Esto implica que:
\begin{la}
    \item Para que haya rendimientos decrecientes a escala, todas las elasticidades-producto de cada factor habrán de ser menores que 1 (si las PMg de todos los factores son positivas);
    \item Para que haya rendimientos constantes a escala, las elasticidades-producto de cada factor no podrán ser superiores a 1;
    \item Mientras que, para que haya rendimientos crecientes a escala, pueden darse tanto con elasticidades-producto mayores o menores que uno (en el caso de varios factores).
\end{la}

\imagenfit{RendimientosEscala_Productividad.png}{Rendimientos a escala y productividadde un factor}{Apuntes Víctor Gutiérrez. Tema 3.A.11}

\subsubsection{Funciones de producción homogéneas}
Una función es homogénea de un cierto grado, si al variar todos sus argumentos (factores, en nuestro caso) multiplicándolos por un escalar cualquiera ($\lambda$), el valor de la función (el producto, en nuestro caso) se ve multiplicado por un valor (sea $\lambda'$). Esto es:

\eqblock{
\begin{aligned}
    &f(\lambda K,\lambda L)=\lambda'f(K,L)\\[2em]
    &Q=\frac{\partial Q}{\partial L}\;L+\frac{\partial Q}{\partial K}K\;\Longrightarrow\;
    \begin{aligned}
        &L^{-1}\;\Longrightarrow\quad
        \left\{\begin{aligned}
            &\frac{Q}{L}=\frac{\partial Q}{\partial L}+\frac{\partial Q}{\partial K}\frac{K}{L}\\[0.5em]
            &PMe_L=PMg_L+PMg_K\frac{K}{L}
        \end{aligned}\right.\\[0.5em]
        &K^{-1}\;\Longrightarrow\quad
        \left\{\begin{aligned}
            &\frac{Q}{L}=\frac{\partial Q}{\partial K}\frac{L}{K}+\frac{\partial Q}{\partial K}\\[0.5em]
            &PMe_K=PMg_L\frac{L}{K}+PMg_K
        \end{aligned}\right.\\[0.5em]
    \end{aligned}
\end{aligned}
}{Funciones de Producción homogeneas. Teorema de Euler. Leyes de producción: Largo plazo}

En este caso, ver los rendimientos a escala de la función de producción es muy sencillo, puesto que el grado de homogeneidad recoge precisamente la noción de rendimientos. Si el grado de homogeneidad es mayor a uno, la función presentará RCrE, si el grado de homogeneidad es igual a uno, la función tendrá RCE y, por último, si el grado de homogeneidad está (0,1), entonces presenta RDE.
En este caso, si se aplica el Teorema de Euler (homogeneidad de funciones), obtendremos que:
 
Si la $PMe_L=PMg_L$, la $PMg_K$ se iguala a cero. Análogamente, si la $PMe_K=PMg_K$, entonces $PMg_L$ se iguala a cero. De ahí que si la función de producción es homogenea, en el punto donde $PMe_L=PMg_L$, la $PMg_K$ (factor inelástico en este caso) se hace cero. Esta da un sentido particularmente claro a la FASE II, ya que está queda acotada por las PMg’s nulas de ambos factores. Es decir, en una función homogenea, los límites de la FASE II quedan determinados justamente por los puntos donde los productos marginales de los factores se anulan, ya que (indpendientemente del factor elástico):
\begin{lnun}
    \item Cuando $PMe_L=PMg_L$, se cumple que $PMg_K=0$; entonces:
    \begin{la}
        \item Si $PMe_L > PMg_L \;\Longrightarrow\; PMg_K > 0$
        \item Si $PMe_L < PMg_L \;\Longrightarrow\; PMg_K < 0$
    \end{la}
    \item Cuando $PMe_K=PMg_K$, se cumple que $PMg_L=0$; entonces
    \begin{la}
        \item Si $PMeK > PMgK ⟹ PMgL > 0$
        \item Si $PMe_K < PMg_K\;\Longrightarrow\;PMg_L < 0$
    \end{la}
\end{lnun}

\imagenfit{FASES_CASSEL.png}{Diagrama de Fases de CASSEL}{Apuntes Víctor Gutiérrez. Tema 3.A.11}
	
Entonces, dentro de la FASE II, ambos productos marginales son positivos, y fuera de ella, uno u otro es cero o negativo. De esta forma, si la PMe de un factor fuese inferior a su PMg, la PMg del otro tendría que ser negativa. Viceversa, para que sea positiva, la PMe del primer factor debeía ser superior a su PMg. Por ello se contrarán factores únicamente en el tramo en que $PMe \ge PMg \ge 0$.\footnote{%
    La aplicación de esto tiene enormes implicaciones en materia de distribución de la renta entre factores en el caso de funciones de producción agregadas. Si estamos en competencia perfecta: $PMg_L=w$ y la $PMg_K=r$, de donde se cumple que la misma ecuación inicial $X=PMg_LL+PMg_KK=wL+rK$. Por lo que se cumple el Teorema del agotamiento del producto al ir todo él a salarios ($wL$) o a beneficios ($rK$)
} 
Una condición semejante no existe para el caso de funciones no homogéneas, en este caso pasaría a ser contingente.

\subsection{Muy largo plazo: Progreso técnico}
A modo de extensión, es interesante incluir un horizonte temporal que podríamos denominar muy largo plazo,\footnote{%
    Para profundizar en el concepto de progreso técnico: HIWELL JONES, \textit{Introducción a las teorías modernas del crecimiento económico}, Capítulo VII \textit{Las acepciones sencillas de progreso técnico}, pág. 189-222
} en el que asumimos que la morfología del conjunto de posibilidades de producción puede cambiar. Dada la naturaleza de la cuestión, es algo que se estudia esencialmente en las teorías del crecimiento económico, por tanto, en el campo de la macroeconomía. Es importante por tanto entender que, mientras que hasta ahora hemos tratado la tecnología de una empresa individual, lo que sigue se refiere a la tecnología agregada de la economá.

El progreso técnico supone nuevos y mejores métodos de producción, o una mejora en la eficiencia de los métodos ya existentes. Esto se puede ver reflejado de dos maneras:
\begin{la}
    \item Innovaciones en el proceso productivo, que permiten que puedan producirse más bienes utilizando las mismas cantidades de factores (o que se pueda obtener la misma cantidad de bien con cantidades menores de uno o más factores); y,
    \item Innovaciones en el producto, que permiten que los productos existentes mejoren su calidad o que se produzcan bienes completamente nuevos.
\end{la}

El progreso técnico puede representarse gráficamente como un desplazamiento hacia arriba de la función de producción que se traduce en un desplazamiento hacia adentro de las curvas isocuantas:

\imagenfit{ISOCUANTAS_LP.png}{Progreso técnico}{Apuntes Víctor Gutiérrez. Tema 3.A.11}
 
Consideremos una función de producción agregada $\hat F$ y definiremos los siguientes tipos de progreso técnico neutral.

\subsubsection{Progreso técnico: HICKS}
Progreso técnico neutral en sentido de HICKS.

$$\hat F(A_t, z_{K_t}, z_{L_t})=A_t\;\hat F(z_{K_t}, z_{L_t})$$

En el progreso técnico neutral de HICKS, el capital y el trabajo son afectados por el progreso técnico. La cantidad de factores utilizados disminuye. Aumenta la eficiencia y la productividad de todos los inputs utilizados.

\imagenfit{MLP_HICKS.png}{Progreso tecnológico en el sentido de HICKS. Análisis del desplazamiento delas isocuantas}{Apuntes Víctor Gutiérrez. Tema 3.A.11}

\subsubsection{Progreso técnico: SOLOW}
Progreso técnico neutral en sentido de SOLOW. Intensificador del capital.

$$\hat F(A_t, z_{K_t}, z_{L_t})=F(A_t\cdot z_{K_t}, z_{L_t})$$
 
En el progreso técnico neutral de SOLOW, la función de producción define al capital afectado por el progreso técnico (intensificado). El capital es medido en términos de unidades de eficiencia y la relación producto capital no es constante. El capital es más eficiente.

\imagenfit{MLP_SOLOW.png}{Progreso tecnológico en el sentido de SOLOW. Análisis del desplazamiento delas isocuantas}{Apuntes Víctor Gutiérrez. Tema 3.A.11}

\subsubsection{Progreso técnico: HARROD}
Progreso técnico neutral en sentido de HARROD. Intensificador del trabajo. 

$$\hat F(A_t, z_{K_t}, z_{L_t})=F(z_{K_t}, A_t\cdot z_{L_t})$$

En el progreso técnico neutral de HARROD, la función de producción es modificada por un trabajo que aprende. 

\imagenfit{MLP_HARROD.png}{Progreso tecnológico en el sentido de HARROD. Análisis del desplazamiento delas isocuantas}{Apuntes Víctor Gutiérrez. Tema 3.A.11}
 








\clearpage
\section{Conclusión} 

\subsection{Extensiones y relación con otras partes del Temario}


\subsection{Opinión}



\subsubsection{Idea final (Salida o cierra)}


\clearpage
\pagenumbering{gobble}
\MyBibliography
\MyBibItem{Autor}{Título}{Año}{DOI -- LINK}




\anexo{\textbf{Preguntas Test}}
%\input{\TemaCode/questions.tex}






\end{document}